% AY2025 力学 解答例
% 体裁・粒度は 参考/ のPDFに揃える（行間を埋め、最終解答は \ans{} で boxed）。
% 解答・数値・問題は本年度過去問から自力で導いた（東大版は体裁の手本のみ）。
\documentclass[11pt,a4paper]{article}
\usepackage{../../../00_common/preamble}

\usetikzlibrary{decorations.pathmorphing}

\title{力学 AY2025 解答例（専門基礎科目）}
\author{}
\date{}

\begin{document}
\maketitle

% ==================================================================
\section*{第1問〔I〕}

質量 $m_1$ の質点 1，質量 $m_2$ の質点 2 の間に，作用反作用の力のみが働く。
原点 O から測った各質点の位置ベクトルを $\bm{r}_1$，$\bm{r}_2$ とする。
質点 1 が質点 2 から受ける力を $\bm{F}_{12}$，質点 2 が質点 1 から受ける力を
$\bm{F}_{21}$ とすると，作用反作用の法則により
\[
  \bm{F}_{12}=-\bm{F}_{21},\qquad |\bm{F}_{12}|=|\bm{F}_{21}|=F.
\]
外力は働かないので，各質点の運動方程式は
\[
  m_1\ddot{\bm{r}}_1=\bm{F}_{12},\qquad
  m_2\ddot{\bm{r}}_2=\bm{F}_{21}=-\bm{F}_{12}.
\]

\subsection*{(1) 運動量保存則}
系の全運動量を $\bm{P}=m_1\dot{\bm{r}}_1+m_2\dot{\bm{r}}_2$ とおく。
これを時刻 $t$ で微分し，上の二つの運動方程式を代入すると
\begin{align*}
  \frac{d\bm{P}}{dt}
  &=m_1\ddot{\bm{r}}_1+m_2\ddot{\bm{r}}_2
   =\bm{F}_{12}+\bm{F}_{21}
   =\bm{F}_{12}+(-\bm{F}_{12})
   =\bm{0}.
\end{align*}
よって全運動量の時間変化は零であり，$\bm{P}$ は時間によらず一定に保たれる。
\[
  \ans{\;\bm{P}=m_1\dot{\bm{r}}_1+m_2\dot{\bm{r}}_2=\text{一定}
  \quad\Bigl(\because\ \frac{d\bm{P}}{dt}=\bm{0}\Bigr)\;}
\]
すなわち，内力（作用反作用の対）は系の全運動量を変化させず，運動量保存則が成り立つ。

\subsection*{(2) 一方の質点から見た他方の運動方程式}
質点 1 の上に立つ観測者が質点 2 の運動を観測する場合を考える
（質点 1 と 2 を入れ替えても同形になる）。観測される相対位置ベクトルを
\[
  \bm{r}\equiv\bm{r}_2-\bm{r}_1
\]
とおく。二つの運動方程式をそれぞれの質量で割って差をとると
\begin{align*}
  \ddot{\bm{r}}=\ddot{\bm{r}}_2-\ddot{\bm{r}}_1
  &=\frac{\bm{F}_{21}}{m_2}-\frac{\bm{F}_{12}}{m_1}
   =\frac{\bm{F}_{21}}{m_2}-\frac{-\bm{F}_{21}}{m_1}
   =\left(\frac{1}{m_1}+\frac{1}{m_2}\right)\bm{F}_{21}.
\end{align*}
ここで換算質量 $\mu$ を $\dfrac{1}{\mu}=\dfrac{1}{m_1}+\dfrac{1}{m_2}$，
すなわち $\mu=\dfrac{m_1 m_2}{m_1+m_2}$ で定めると，
\[
  \ans{\;\mu\,\ddot{\bm{r}}=\bm{F}_{21},\qquad
  \mu=\dfrac{m_1 m_2}{m_1+m_2}\;}
\]
となる。質点 1 上の（一般には非慣性系の）観測者から見ると，質点 2 はあたかも
質量が換算質量 $\mu$ の 1 質点であるかのように，相手から受ける力 $\bm{F}_{21}$
のもとで運動する。観測者が質点 2 の側に立つ場合も，相対座標を
$\bm{r}_1-\bm{r}_2$ として同様に $\mu(\ddot{\bm{r}}_1-\ddot{\bm{r}}_2)=\bm{F}_{12}$
が得られる。

\medskip
検算: (2) で $m_2\to\infty$（相手が無限に重い）とすると
$\mu\to m_1$ かつ $\bm{r}\to-\bm{r}_1$ となり，$m_1\ddot{\bm{r}}_1=\bm{F}_{12}$
という質点 1 単体の運動方程式に帰着する。固定された重い質点を基準にした観測と
一致し，妥当。また (1) の運動量保存と (2) は両立する（重心は等速直線運動，
相対運動は $\mu$ で記述される）。✓

% ==================================================================
\clearpage
\section*{第2問〔II〕}

半径 $a$，質量 $M$，重心まわりの慣性モーメント $I=\dfrac{1}{2}Ma^2$ の均一円盤を，
水平とのなす角 $\theta$ の粗い斜面上に静かに置くと，滑らずに回転しながら斜面を
下る。重力加速度の大きさを $g$，静止摩擦係数を $\mu$ とする。

斜面下向きを正にとり，重心の加速度の大きさを $\alpha_G$，角加速度の大きさを
$\dot\omega$ とする。円盤に働く力は，斜面に沿う重力成分 $Mg\sin\theta$（下向き），
斜面に垂直な抗力 $N$，および接点に働く摩擦力 $f$（滑り出しを妨げる向き，
すなわち斜面上向き）である。

\subsection*{(1) 摩擦力の大きさ}
斜面に垂直な方向のつり合いより
\[
  N=Mg\cos\theta.
\]
斜面に沿う方向の重心の並進運動方程式と，重心まわりの回転の運動方程式は
\begin{align*}
  &\text{並進}:\quad M\alpha_G = Mg\sin\theta - f,\\
  &\text{回転}:\quad I\dot\omega = f\,a
    \quad\Longrightarrow\quad \frac{1}{2}Ma^2\,\dot\omega=f\,a.
\end{align*}
滑らずに転がる（転がり条件）から，接点速度が零で $\alpha_G=a\,\dot\omega$，
すなわち $\dot\omega=\dfrac{\alpha_G}{a}$。これを回転の式へ代入すると
\[
  \frac{1}{2}Ma^2\cdot\frac{\alpha_G}{a}=f\,a
  \quad\Longrightarrow\quad
  f=\frac{1}{2}M\alpha_G.
\]
これと並進の式 $M\alpha_G=Mg\sin\theta-f$（すなわち $M\alpha_G=2f$）から
$\alpha_G$ を消去すると
\[
  2f=Mg\sin\theta-f
  \quad\Longrightarrow\quad
  3f=Mg\sin\theta.
\]
\[
  \ans{\;f=\dfrac{1}{3}Mg\sin\theta\;}
\]

\subsection*{(2) 滑らずに転がる条件}
摩擦力は静止摩擦の範囲になければならない。すなわち
\[
  f\le \mu N=\mu Mg\cos\theta.
\]
(1) の $f=\dfrac{1}{3}Mg\sin\theta$ を代入して
\[
  \frac{1}{3}Mg\sin\theta\le \mu Mg\cos\theta
  \quad\Longrightarrow\quad
  \tan\theta\le 3\mu.
\]
\[
  \ans{\;\tan\theta\le 3\mu
  \quad\Bigl(\text{すなわち}\ \mu\ge\dfrac{\tan\theta}{3}\Bigr)\;}
\]

\subsection*{(3) 重心の加速度の大きさ}
並進の式 $M\alpha_G=2f$ に (1) の $f=\dfrac{1}{3}Mg\sin\theta$ を代入すると
\[
  M\alpha_G=2\cdot\frac{1}{3}Mg\sin\theta
  \quad\Longrightarrow\quad
  \alpha_G=\frac{2}{3}g\sin\theta.
\]
\[
  \ans{\;\alpha_G=\dfrac{2}{3}g\sin\theta\;}
\]

\medskip
検算: エネルギーから確認する。斜面に沿って重心が距離 $\ell$ だけ下ると，
転がり条件 $v_G=a\omega$ のもとでエネルギー保存は
\[
  Mg\ell\sin\theta=\frac{1}{2}Mv_G^2+\frac{1}{2}I\omega^2
  =\frac{1}{2}Mv_G^2+\frac{1}{2}\cdot\frac{1}{2}Ma^2\cdot\frac{v_G^2}{a^2}
  =\frac{3}{4}Mv_G^2.
\]
ゆえに $v_G^2=\dfrac{4}{3}g\ell\sin\theta$。等加速度運動 $v_G^2=2\alpha_G\ell$ と
比べると $\alpha_G=\dfrac{2}{3}g\sin\theta$ となり，(3) と一致する。
また摩擦が無い場合（$f=0$）の純粋な滑り落ち $g\sin\theta$ より小さく，
回転にエネルギーが配分される分だけ加速度が減ることとも整合する。✓

\end{document}
